\documentclass[CJK,13pt]{beamer}
\usepackage{CJKutf8}
\usepackage{beamerthemesplit}
\usetheme{Warsaw}
\useoutertheme[footline=authortitle]{miniframes}
\usepackage{amsmath}
\usepackage{amssymb}
\usepackage{graphicx}
\usepackage{eufrak}
\usepackage{color}
\usepackage{slashed}
\usepackage{simplewick}
\usepackage{tikz}
\usepackage{tcolorbox}
\usepackage{ulem}
%%figures
\def\lfig#1#2{\includegraphics[width=#1 in]{#2}}
\def\tfig#1#2{\includegraphics[height=#1 in]{#2}}
\def\addfig#1#2{\begin{center}\includegraphics[width=#1 in]{#2}\end{center}}
\def\question{\includegraphics[width=0.3in]{why.jpg}\,}
\def\answer{\includegraphics[width=0.3in]{baozou_haha.png}\,}
\def\wulian{\includegraphics[width=0.18in]{emoji_wulian.jpg}}
\def\bigwulian{\includegraphics[width=0.35in]{emoji_wulian.jpg}}
\def\bye{\includegraphics[width=0.18in]{emoji_bye.jpg}}
\def\bigbye{\includegraphics[width=0.35in]{emoji_bye.jpg}}
\def\huaixiao{\includegraphics[width=0.18in]{emoji_huaixiao.jpg}}
\def\bighuaixiao{\includegraphics[width=0.35in]{emoji_huaixiao.jpg}}
\def\jianxiao{\includegraphics[width=0.18in]{emoji_jianxiao.jpg}}
\def\bigjianxiao{\includegraphics[width=0.35in]{emoji_jianxiao.jpg}}
\def\haoqi{\includegraphics[width=0.18in]{emoji_haoqi.jpg}}
%% colors
\def\blacktext#1{{\color{black}#1}}
\def\bluetext#1{{\color{blue}#1}}
\def\redtext#1{{\color{red}#1}}
\def\darkbluetext#1{{\color[rgb]{0,0.2,0.6}#1}}
\def\skybluetext#1{{\color[rgb]{0.2,0.7,1.}#1}}
\def\cyantext#1{{\color[rgb]{0.,0.5,0.5}#1}}
\def\greentext#1{{\color[rgb]{0,0.7,0.1}#1}}
\def\darkgray{\color[rgb]{0.2,0.2,0.2}}
\def\lightgray{\color[rgb]{0.6,0.6,0.6}}
\def\gray{\color[rgb]{0.4,0.4,0.4}}
\def\blue{\color{blue}}
\def\red{\color{red}}
\def\orange{\color[rgb]{1.,0.8,0.}}
\def\green{\color{green}}
\def\darkgreen{\color[rgb]{0,0.4,0.1}}
\def\darkblue{\color[rgb]{0,0.2,0.6}}
\def\skyblue{\color[rgb]{0.2,0.7,1.}}
%%control
\def\ulim#1{\underset{#1}{\lim}}
\def\heaviside{\,\mathrm{h}\,}
\def\bral{\left(\begin{array}{l}}
\def\brar{\end{array}\right)}
\def\brall{\left(\begin{array}{ll}}
\def\brarr{\end{array}\right)}
\def\bralll{\left(\begin{array}{lll}}
\def\brarrr{\end{array}\right)}
\def\branchl{\left\{\begin{array}{l}}
\def\branchr{\end{array}\right.}
\def\branchll{\left\{\begin{array}{ll}}
\def\branchrr{\end{array}\right.}
\def\branchlll{\left\{\begin{array}{lll}}
\def\branchrrr{\end{array}\right.}
\def\sfgamma#1{\,\Gamma\left( #1 \right)\,}
\def\be{\begin{equation}}
\def\ee{\nonumber\end{equation}}
\def\bea{\begin{eqnarray}}
\def\eea{\nonumber\end{eqnarray}}
\def\bch{\begin{CJK}{UTF8}{gbsn}}
\def\ech{\end{CJK}}
\def\bitem{\begin{itemize}}
\def\eitem{\end{itemize}}
\def\bcenter{\begin{center}}
\def\ecenter{\end{center}}
\def\bex{\begin{minipage}{0.2\textwidth}\includegraphics[width=0.6in]{jugelizi.png}\end{minipage}\begin{minipage}{0.76\textwidth}}
\def\eex{\end{minipage}}
\def\chtitle#1{\frametitle{\bch#1\ech}}
\def\bmat#1{\left(\begin{array}{#1}}
\def\emat{\end{array}\right)}
\def\bcase#1{\left\{\begin{array}{#1}}
\def\ecase{\end{array}\right.}
\def\bmini#1{\begin{minipage}{#1\textwidth}}
\def\emini{\end{minipage}}
\def\tbox#1{\begin{tcolorbox}#1\end{tcolorbox}}
\def\pfrac#1#2#3{\left(\frac{\partial #1}{\partial #2}\right)_{#3}}
\def\res#1#2{\,\mathrm{res}\,#1\left(#2\right)\,}
\def\newt#1#2{\left(\begin{array}{c}#1\\ #2\end{array}\right)}
\def\reof#1{\,\mathrm{Re}{\left(#1\right)\,}}
\def\imof#1{\,\mathrm{Im}{\left(#1\right)\,}}
\def\Arg#1{\,\mathrm{Arg}\,#1\,}
%%symbols
\def\intfull{\int_{-\infty}^\infty}
\def\inthalf{\int_0^\infty}
\def\bropt{\,(\ \ \ )}
\def\sech{\mathrm{sech}\,}
\def\csch{\mathrm{csch}\,}
\def\sone{$\star$}
\def\stwo{$\star\star$}
\def\sthree{$\star\star\star$}
\def\sfour{$\star\star\star\star$}
\def\sfive{$\star\star\star\star\star$}
\def\rint{{\int_\leftrightarrow}}
\def\roint{{\oint_\leftrightarrow}}
\def\stdHf{{\textit{\r H}_f}}
\def\deltaH{{\Delta \textit{\r H}}}
\def\ii{{\dot{\imath}}}
\def\skipline{{\vskip0.1in}}
\def\skiplines{{\vskip0.2in}}
\def\lagr{{\mathcal{L}}}
\def\hamil{{\mathcal{H}}}
\def\vecv{{\mathbf{v}}}
\def\vecf{{\mathbf{f}}}
\def\vecg{{\mathbf{g}}}
\def\vecupsilon{{\mathbf{\upsilon}}}
\def\vecu{{\mathbf{u}}}
\def\vecj{{\mathbf{j}}}
\def\vecx{{\mathbf{x}}}
\def\vecy{{\mathbf{y}}}
\def\vecz{{\mathbf{z}}}
\def\veck{{\mathbf{k}}}
\def\vecp{{\mathbf{p}}}
\def\vecn{{\mathbf{n}}}
\def\vecA{{\mathbf{A}}}
\def\vecP{{\mathbf{P}}}
\def\vecsigma{{\mathbf{\sigma}}}
\def\hatJn{{\hat{J_\vecn}}}
\def\hatJx{{\hat{J_x}}}
\def\hatJy{{\hat{J_y}}}
\def\hatJz{{\hat{J_z}}}
\def\hatj#1{\hat{J_{#1}}}
\def\hatphi{{\hat{\phi}}}
\def\hatq{{\hat{q}}}
\def\hatpi{{\hat{\pi}}}
\def\vel{\upsilon}
\def\Dint{{\mathcal{D}}}
\def\adag{{\hat{a}^\dagger}}
\def\bdag{{\hat{b}^\dagger}}
\def\cdag{{\hat{c}^\dagger}}
\def\ddag{{\hat{d}^\dagger}}
\def\hata{{\hat{a}}}
\def\hatb{{\hat{b}}}
\def\hatc{{\hat{c}}}
\def\hatd{{\hat{d}}}
\def\hatD{{\,\hat{D}}}
\def\hatN{{\hat{N}}}
\def\hatH{{\hat{H}}}
\def\hatp{{\hat{p}}}
\def\Fup{{F^{\mu\nu}}}
\def\Fdown{{F_{\mu\nu}}}
\def\newl{\nonumber \\}
\def\vece{\mathrm{e}}
\def\calM{{\mathcal{M}}}
\def\calT{{\mathcal{T}}}
\def\calR{{\mathcal{R}}}
\def\barpsi{\bar{\psi}}
\def\baru{\bar{u}}
\def\barv{\bar{\upsilon}}
\def\qeq{\stackrel{?}{=}}
\def\ftf{\stackrel{\mathcal{FT}}{\Longrightarrow}}
\def\ftb{\stackrel{\mathcal{FT}}{\Longleftarrow}}
\def\ftfb{\stackrel{\mathcal{FT}}{\Longleftrightarrow}}
\def\ltf{\stackrel{\mathcal{LT}}{\Longrightarrow}}
\def\ltb{\stackrel{\mathcal{LT}}{\Longleftarrow}}
\def\ltfb{\stackrel{\mathcal{LT}}{\Longleftrightarrow}}
\def\torder#1{\mathcal{T}\left(#1\right)}
\def\rorder#1{\mathcal{R}\left(#1\right)}
\def\contr#1#2{\contraction{}{#1}{}{#2}#1#2}
\def\trof#1{\mathrm{Tr}\left(#1\right)}
\def\trace{\mathrm{Tr}}
\def\comm#1{\ \ \ \left(\mathrm{used}\ #1\right)}
\def\tcomm#1{\ \ \ (\text{#1})}
\def\slp{\slashed{p}}
\def\slk{\slashed{k}}
\def\calp{{\mathfrak{p}}}
\def\veccalp{\mathbf{\mathfrak{p}}}
\def\Tthree{T_{\tiny \textcircled{3}}}
\def\pthree{p_{\tiny \textcircled{3}}}
\def\dbar{{\,\mathchar'26\mkern-12mu d}}
\def\erf{\mathrm{erf}}
\def\const{\mathrm{const.}}
\def\pheat{\pfrac p{\ln T}V}
\def\vheat{\pfrac V{\ln T}p}
%%units
\def\fdeg{{^\circ \mathrm{F}}}
\def\cdeg{^\circ \mathrm{C}}
\def\atm{\,\mathrm{atm}}
\def\angstrom{\,\text{\AA}}
\def\SIL{\,\mathrm{L}}
\def\SIkm{\,\mathrm{km}}
\def\SIyr{\,\mathrm{yr}}
\def\SIGyr{\,\mathrm{Gyr}}
\def\SIV{\,\mathrm{V}}
\def\SImV{\,\mathrm{mV}}
\def\SIeV{\,\mathrm{eV}}
\def\SIkeV{\,\mathrm{keV}}
\def\SIMeV{\,\mathrm{MeV}}
\def\SIGeV{\,\mathrm{GeV}}
\def\SIcal{\,\mathrm{cal}}
\def\SIkcal{\,\mathrm{kcal}}
\def\SImol{\,\mathrm{mol}}
\def\SIN{\,\mathrm{N}}
\def\SIHz{\,\mathrm{Hz}}
\def\SIm{\,\mathrm{m}}
\def\SIcm{\,\mathrm{cm}}
\def\SIfm{\,\mathrm{fm}}
\def\SImm{\,\mathrm{mm}}
\def\SInm{\,\mathrm{nm}}
\def\SImum{\,\mathrm{\mu m}}
\def\SIJ{\,\mathrm{J}}
\def\SIW{\,\mathrm{W}}
\def\SIkJ{\,\mathrm{kJ}}
\def\SIs{\,\mathrm{s}}
\def\SIkg{\,\mathrm{kg}}
\def\SIg{\,\mathrm{g}}
\def\SIK{\,\mathrm{K}}
\def\SImmHg{\,\mathrm{mmHg}}
\def\SIPa{\,\mathrm{Pa}}
%page
\def\secpage#1#2{\begin{frame}\bcenter{\bf \Huge #1} \ecenter \skipline \tbox{\bcenter #2 \ecenter}\end{frame}}
\def\append#1#2{\secpage{附录#1}{#2}}
\def\thinka#1{\begin{frame}\frametitle{思考题}\bcenter\lfig{0.4}{think.jpg}  \skipline #1 \ecenter\end{frame}}
\def\thinkb#1{\begin{frame}\frametitle{思考题}\bcenter\lfig{0.5}{think1.jpg}  \skipline #1 \ecenter\end{frame}}
\def\thinkc#1{\begin{frame}\frametitle{思考题}\bcenter\lfig{0.5}{think2.jpg}  \skipline #1 \ecenter\end{frame}}
\def\thinkd#1{\begin{frame}\frametitle{思考题}\bcenter\lfig{0.5}{think3.jpg}  \skipline #1 \ecenter\end{frame}}
\def\thinke#1{\begin{frame}\frametitle{思考题}\bcenter\lfig{0.5}{think4.jpg}  \skipline #1 \ecenter\end{frame}}

\def\courseurl{http://zhiqihuang.top/cosm}

\def\tpage#1#2{
\title{COSM \S{#1}  #2}
  \author{Zhiqi Huang}

\begin{frame}
\begin{center}
{\bf \Large Cosmology}

{\vskip 0.1in}

{\Large \S #1 #2}

{\vskip 0.2in}

{Lecturer: 黄志琦}

\vskip 0.2in

\courseurl

\end{center}
\end{frame}
}


  \date{}
  \begin{document}
  \bch
  \tpage{12}{Matter Power Spectrum and CMB Power Spectrum}

  \secpage{下面我们介绍如何用前面介绍的微扰计算结果来计算宇宙学观测量}{\addfig{1}{qidai.jpg}}

  
  \section{Linear Matter Power Spectrum}

  \secpage{冷物质线性扰动功率谱}{$$P_m(k; z) = \frac{8\pi^2a^2k}{9H_0^2\Omega_m^2} \lvert\Phi(k;z)\rvert^2 \Delta^2(k).$$}
  
  \begin{frame}
    \frametitle{冷物质的密度扰动}
    利用泊松方程，可以从牛顿引力势算出“冷物质密度扰动” $\delta_m$：
    $$ 4\pi G\bar{\rho}_m \delta_m = -\frac{k^2}{a^2}\Phi $$
    注意密度扰动是依赖于规范的量，这里的 $\delta_m$ 和牛顿规范下的 $\delta_c, \delta_b$ 并不完全一致。它们仅在晚期宇宙、远小于视界尺度（$k\gg aH$） 时一致。在忽略辐射和中微子的情况下，$\delta_m$ 和物质共动同步规范下的冷物质密度扰动一致。
  \end{frame}

  \begin{frame}
    \frametitle{冷物质密度线性扰动功率谱}
    利用 $\bar{\rho}_m = \frac{3H_0^2}{8\pi G}\Omega_m a^{-3}$，可以写出冷物质密度扰动的功率谱为
    $$P_m(k; z) = \frac{4a^2k^4}{9H_0^2\Omega_m^2} \lvert\Phi(k;z)\rvert^2 P_\zeta(k)$$
    我们代码算出来的 $\Phi(k; z)$ 实际上是隐去了原初扰动初始随机变量 $\zeta(\vec{k})$ 的因子，所以上面补回了 $\zeta$ 的功率谱。

    用 $\zeta$ 的无量纲功率谱  $\Delta^2(k)=\frac{k^3}{2\pi^2}P_\zeta(k)$ 表示就是
    $$P_m(k; z) = \frac{8\pi^2a^2k}{9H_0^2\Omega_m^2} \lvert\Phi(k;z)\rvert^2 \Delta^2(k).$$
  \end{frame}


  \begin{frame}
    \frametitle{废话不多，直接上代码}
    \url{http://zhiqihuang.top/codes/matterpower.py}

    \skipline
    
    因为可能需要输出很大的 $k$ 的 $P_m(k)$，代码里对大的 $k$ 有特殊的处理方式：主要是把晚期宇宙的辐射和中微子的多极展开中的 $\ell>2$ 项进行人为压缩。物理上，由于晚期宇宙辐射和中微子占比低，且在小尺度($k\gg aH$)上很均匀，所以这个近似应该OK（数值上，你也可以通过改变压缩量来检测事实是否如此）。

    \skipline
    
    另外，要注意这里计算的是线性扰动功率谱。实际上的物质功率谱在 $k\gtrsim 0.1\mathrm{Mpc}^{-1}$ 逐渐会有非线性的修正。非线性的物质扰动功率谱一般是用大尺度结构数值模拟和一些半解析方法进行拟合得到。
  \end{frame}

  \begin{frame}
    \frametitle{$\sigma_8$ 参数}
    宇宙学的 $\sigma_8$ 参数定义为红移零处，半径为 $R = 8h^{-1}\mathrm{Mpc}$ 的球内的物质的密度扰动平均值的宇宙学标准差。也就是说，假想实现很多宇宙，在每个宇宙内取个这样的球，算出球内平均 $\overline{\delta_m}$，然后计算所有宇宙的 $ \overline{\delta_m}$ 标准差（即方差的平方根)。

      具体来说，
      $$\sigma_8^2 = \frac{9}{2\pi^2}\int_0^\infty k^2P_m(k)\left(\frac{\sin(kR)-kR\cos{kR}}{(kR)^3}\right)^2 dk$$

      你能用卷积定理和傅立叶变换保内积的定理给出上面公式的证明吗？


      \skipline

      \notes{宇宙学实测的$\sigma_8$ 大致在 $0.8$ 左右，各个观测的结果之间小有差异，但都稳定分布在 $0.7$ 和 $0.9$ 之间。}
      
  \end{frame}

  
  
  \section{CMB power spectrum}
  
  \secpage{CMB power spectrum}{加法公式又来了……}

  \begin{frame}
    \frametitle{宇宙微波背景辐射的光强扰动}
    由于再电离的光深不大，原初宇宙光子可以在地球附近被观测到。由于目前宇宙原初光子的温度 $T_{\rm CMB}=2.726\SIK$ 在微波波段，这些光子就被称为宇宙微波背景辐射（Cosmic Microwave Background），简称CMB。当然，实际的观测并非如此简单，我们至少需要避开银河系和河外星系的星光的干扰。如果是地面上的观测，我们还要考虑大气干扰等各种因素。

    \skipline

    我们可以在两个方向 $\vec{n}_1$ 和 $\vec{n}_2$ 观测光子的分布函数扰动（偏离黑体谱）。一般我们只关心对 $q^3dq$ 积分后的光强扰动。它可以由傅立叶逆变换给出：    
    $$\frac{\Delta I}{I}(\tau, \vec{x}, \vec{n}) = \frac{1}{(2\pi)^3}\int d^3\vec{k} \sum_{\ell}(-i)^\ell F_{\gamma\ell}(\tau, k)\zeta(\vec{k}) P_{\ell}(\vec{n}\cdot\hat{k})e^{i\vec{k}\cdot\vec{x}}$$

  \end{frame}

  \begin{frame}
    \frametitle{宇宙微波背景辐射温度扰动}    
    我们只能在 $\tau = \tau_0$ (红移零处)观测CMB，另外根据宇宙学原理，也可以不妨取地球附近 $\vec{x} = 0$ (这一步不是必须的，即使不嫌麻烦留着 $\vec{x}$， 最后也会消掉)。另外由于 $I\propto T^4$，所以观测到的
    $$\frac{\Delta T}{T}(\vec{n}) =\frac{1}{4} \frac{1}{(2\pi)^3}\int d^3\vec{k} \sum_{\ell}(-i)^\ell F_{\gamma\ell}(\tau_0, k)\zeta(\vec{k}) P_{\ell}(\vec{n}\cdot\hat{k}) $$
    由于存在随机变量 $\zeta(\vec{k})$，我们无法预言任何一个方向的 $\frac{\Delta T}{T}(\vec{n})$。所以，我们的兴趣转移到了两个不同方向的温度扰动的统计关联：
  \end{frame}

  \begin{frame}
    \frametitle{宇宙微波背景辐射的角关联函数}
    考虑 $\vec{n_1}$, $\vec{n_2}$ 两个方向上 $\frac{\Delta T}{T}$ 的统计关联（这个的意思是，按照暴胀理论，假想随机生成很多个宇宙并对这些宇宙中取平均；我们的宇宙只是随机生成的一个，所以实际观测到的量会偏离平均，这种偏离称为 cosmic variance）：
    \bea
     && \langle \frac{\Delta T}{T}(\vec{n}_1)  \frac{\Delta T}{T}(\vec{n}_2)\rangle \newl
    &=& \frac{1}{16} \frac{1}{(2\pi)^6}\sum_{\ell,\ell'}(-i)^{\ell+\ell'} \times \newl
    && \int d^3\vec{k} \int d^3\vec{k}' F_{\gamma\ell}(\tau_0, k)F_{\gamma\ell'}(\tau_0, k')\langle \zeta(\vec{k})\zeta(\vec{k}')\rangle P_{\ell}(\vec{n_1}\cdot\hat{k}) P_{\ell}(\vec{n_2}\cdot\hat{k}')
    \eea
    
  \end{frame}

  \begin{frame}
    \frametitle{宇宙微波背景辐射的角关联函数（续)}
    代入原初曲率扰动功率谱的定义：
    $$ \langle \zeta(\vec{k})\zeta(\vec{k}')\rangle = (2\pi)^3\delta(\vec{k}+\vec{k}')\frac{2\pi^2}{k^3}\Delta^2(k)$$
      并对 $\vec{k}'$ 积分，得到
    \bea
     && \langle \frac{\Delta T}{T}(\vec{n}_1)  \frac{\Delta T}{T}(\vec{n}_2)\rangle \newl
    &=& \frac{1}{64 \pi}\sum_{\ell,\ell'}(-i)^{\ell+\ell'} \times \newl
    && \int d^3\vec{k}  F_{\gamma\ell}(\tau_0, k)F_{\gamma\ell'}(\tau_0, k)\frac{\Delta^2(k)}{k^3} P_{\ell}(\vec{n_1}\cdot\hat{k}) P_{\ell}(-\vec{n_2}\cdot\hat{k})
    \eea
  \end{frame}

  \begin{frame}
    \frametitle{宇宙微波背景辐射的角关联函数（续)}    
    回忆下球谐函数的加法公式：
    $$ P_\ell(\vec{n}\cdot\hat{k}) = \frac{4\pi}{2\ell+1}\sum_{m=-\ell}^\ell Y_{\ell m}^*(\vec{n})Y_{\ell m}(\hat{k}) .$$
    就有
    \bea
     && \langle \frac{\Delta T}{T}(\vec{n}_1)  \frac{\Delta T}{T}(\vec{n}_2)\rangle \newl
    &=&  \frac{1}{64 \pi}\sum_{\ell,\ell'}\frac{4\pi}{2\ell+1}\frac{4\pi}{2\ell'+1}\sum_{m, m'}(-i)^{\ell+\ell'}Y_{\ell m}^*(\vec{n_1})Y_{\ell' m'}(-\vec{n_2}) \times \newl
    && \int d^3\vec{k}  F_{\gamma\ell}(\tau_0, k)F_{\gamma\ell'}(\tau_0, k)\frac{\Delta^2(k)}{k^3}Y_{\ell m}(\hat{k}) Y_{\ell' m'}^*(\hat{k}) 
    \eea
  \end{frame}
  

    \begin{frame}
    \frametitle{宇宙微波背景辐射的角关联函数（续)}    
      把积分 $d^3\vec{k}$ 写成 $k^2dk d\Omega$，利用球谐函数的正交归一化性质，
      $$\int d^2\Omega Y_{\ell m}(\hat{k}) Y_{\ell' m'}^*(\hat{k}) = \delta_{\ell\ell'}\delta_{mm'}$$
      并对 $\ell', m'$ 求和， 就有
      \bea
      && \langle \frac{\Delta T}{T}(\vec{n}_1)  \frac{\Delta T}{T}(\vec{n}_2)\rangle \newl
      &=& \frac{1}{64 \pi}\sum_{\ell,m}\left[\frac{4\pi}{2\ell+1}\right]^2(-1)^\ell Y_{\ell m}^*(\vec{n_1})Y_{\ell m}(-\vec{n_2}) \times \newl
      && \int_0^\infty  \left[F_{\gamma\ell}(\tau_0, k)\right]^2\Delta^2(k)\frac{dk}{k}
    \eea
  \end{frame}


    \begin{frame}
    \frametitle{宇宙微波背景辐射的角关联函数（续)}    
      利用球谐函数的宇称性质 $(-1)^\ell Y_{\ell m}(-\vec{n}_2) =Y_{\ell m}(\vec{n}_2)$，就有
      \bea
      && \langle \frac{\Delta T}{T}(\vec{n}_1)  \frac{\Delta T}{T}(\vec{n}_2)\rangle \newl
      &=& \frac{1}{64 \pi} \sum_{\ell,m}\left[\frac{4\pi}{2\ell+1}\right]^2 Y_{\ell m}^*(\vec{n_1})Y_{\ell m}(\vec{n_2}) \times \newl
      && \int_0^\infty  \left[F_{\gamma\ell}(\tau_0, k)\right]^2\Delta^2(k)\frac{dk}{k}
    \eea
  \end{frame}
    

    \begin{frame}
    \frametitle{宇宙微波背景辐射的角关联函数（续)}    
      最后，对$\vec{n}_1, \vec{n}_2$ 利用球谐函数加法公式，得到
      \bea
      && \langle \frac{\Delta T}{T}(\vec{n}_1)  \frac{\Delta T}{T}(\vec{n}_2)\rangle \newl
      &=& \frac{1}{16} \sum_{\ell} \frac{1}{2\ell+1} P_\ell(\vec{n}_1\cdot\vec{n}_2) \int_0^\infty  \left[F_{\gamma\ell}(\tau_0, k)\right]^2\Delta^2(k)\frac{dk}{k}
    \eea
  \end{frame}


    \begin{frame}
    \frametitle{宇宙微波背景辐射的角关联函数（完结)}    
      定义CMB功率谱
      $$C_\ell = \frac{\pi}{4(2\ell+1)^2} \int_0^\infty  \left[F_{\gamma\ell}(\tau_0, k)\right]^2\Delta^2(k)\frac{dk}{k}.$$
      就有
      $$ \langle \frac{\Delta T}{T}(\vec{n}_1)  \frac{\Delta T}{T}(\vec{n}_2)\rangle = \sum_{\ell}\frac{2\ell+1}{4\pi}C_\ell P_\ell(\vec{n}_1\cdot\vec{n}_2)$$
  \end{frame}
    


    \begin{frame}
    \frametitle{Line of sight Integral}    
    由于实验可以观测到 $\ell\sim 10^4$ 的CMB功率谱。而我们在解Boltzmann方程时一般会在 $\ell_{\max}$ 约为十几左右做个截断，所以前面推导的公式只在很小的 $\ell$ 具有实际意义。

    \skipline

    实际的计算使用的是 $F_{\gamma\ell}$ 的一个积分表示：
    $$\frac{F_{\gamma\ell}(\tau_0, k)}{2\ell+1} = \int_0^{\tau_0} d\tau e^{-\kappa} S(k,\tau) j_\ell\left(k(\tau_0-\tau)\right)$$
    这里的 $\kappa$ 是积分得到的$\tau$ 和 $\tau_0$ 之间的总光深。源函数 $S(k,\tau)$ 可以用Boltzmann方程给出的 $\ell\le 2$ 的微扰量组合得出（请参考神作 astro-ph/9603033)。
    由于球贝塞尔函数 $j_\ell$ 都可以事先算好或直接用近似算法快速计算，所以这种计算方法可以实现秒算CMB功率谱。
  \end{frame}
    


    \begin{frame}
      \frametitle{回到CAMB和CLASS}
      最后，当你理解了所有计算量的来龙去脉之后，我希望你能回到CAMB或者CLASS而不是依赖于我写的示范代码 —— 毕竟，python太慢，而且中微子是有质量的。

      \skipline

      安装CAMB很简单：

      pip3 install camb

      然后你搜索下camb的示范代码，很快就能学会如何计算各种微扰量、物质功率谱和CMB功率谱。
    \end{frame}
    
    
    \ech
\end{document}




  
